% !TeX root = ../main.tex
% sections/08_conclusion.tex -- Section VIII of the MobiWac 2026 camera-ready (accepted 2026-08-26).
% Base text: the dissertation's Chapter 5 (articles/dissertacao/src/chapters/5_mobiwac/08_conclusion.tex),
% ported 2026-09-06; it already carries the 3.55 / 3.06 margins, the per-axis surfaces and the
% protected sentence P2 ("sharing alone does not explain the outcome"). Camera-ready edits below
% carry a dated comment beside them.
\section{Conclusion}
\label{sec:conclusion}

We tested whether one model can predict both the category and region of a user's next visit. The
central concern was that sharing could help one task while hurting the other. The evidence does not
attribute the outcome to sharing alone. On the category task the input representation moves the
score more than the choice between one model and two: at every one of the six datasets and in every
fold it improves next-category prediction, by $+0.23$ to $+6.29$ macro-F1 points. On the region task, the joint model is also the larger model, and the design does not separate
size from sharing.

On that representation, one model reads both answers from a single saved model. On region it
outperforms the dedicated model at Texas and California and stays within the two-point margin registered before any result was read at the
other four. On category it outperforms the dedicated model at Florida, and the five remaining
differences are equivalent to zero within half a point. The joint model also remains above every
external baseline in Table~\ref{tab:results} at all six datasets, by at least $3.55$ Acc@10 points
over the strongest region reference and $3.06$ macro-F1 points over POI-RGNN on category.

These results do not mean that multi-task learning helps automatically. Whether the exchange of information between the tasks adds anything beyond the representation is
not separated by the evidence here (Section~\ref{sec:discussion}). When the representation preserves the context of
each visit and the architecture keeps a private spatial path where the tasks differ, one model can
predict both what a user will do next and where, in a single forward pass instead of two.
% [2026-09-06, cut C7] Conclusion rewritten from ~510 to ~250 words. Kept: the P2 sentence in its
% repointed form (evidence does not attribute the outcome to sharing alone, with the two measured
% facts), the representation result with its magnitude, the per-axis verdicts with the registered
% surfaces, the external margins 3.55 / 3.06, the "not automatic" closing, the single-model
% property and "instead of two". Dropped: the design-motivation sentences (a fixed vector cannot
% distinguish contexts; the visit vector can include time, nearby places, recent visits), which
% Sections I and III carry; the vocabulary-grouping observation with its two confounders (Sections I
% and VI carry it, and "which is where an auxiliary signal has the most to add" is contradicted by
% the California control); the external-baseline parenthetical (Section VI-C states the protocols).
% No number, verdict or verb changed. No semicolons.
